\documentclass{article}

% --------------------
% Packages
% --------------------
\usepackage[utf8]{inputenc}
\usepackage[T1]{fontenc}
\usepackage[french]{babel}
\usepackage{amsmath,amssymb}
\usepackage{mdframed}
\usepackage{geometry}
\usepackage{pgfplots}
\usepackage{tkz-tab}
\usepackage{bm}
\usepackage{array}
\pgfplotsset{compat=1.18}
\geometry{margin=2cm}

% Pour les coordonnées de vecteurs :
\newcommand*{\Coord}[2]{%
  \ensuremath{
    \begin{pmatrix}
      #1 \\[10pt]
      #2
    \end{pmatrix}}
}

% --------------------
% En-tête / pied (exam)
% --------------------
% \header{}{}{}
% \footer{}{\thepage}{}
\setlength{\parindent}{0pt}

\begin{document}

\begin{center}
    \Large{\textbf{Correction fiche n°3 exercices}}
\end{center}

\textbf{Exercice n$^\circ$1}

On considère les points $A(-2; 2)$, $B(1; 0)$, $C(-1; -3)$.

\begin{enumerate}
    \item Déterminer les coordonnées :

    \begin{enumerate}
        \item Du point $D$ tel que $ABCD$ soit un parallélogramme.

        Avec la propriété du parallélogramme, $ABCD$ est un parallélogramme si : 
        $\bm{\overrightarrow{AB} = \overrightarrow{DC}}$ (et non pas $\overrightarrow{CD}$)
        
        En appliquant la formule on a : $\overrightarrow{AB}\Coord{3}{-2} = 
        \overrightarrow{DC}\Coord{-1-x_D}{-3 - y_D}$

        En résolvant, on obtient les coordonnées de $D$ :
        \[
        D(-4; -1)
        \]

        \item Du centre $O$ de $ABCD$.

        Le centre $O$ est le milieu des diagonales. En appliquand la formule pour trouver le milieu d'un segment,
        on obtient : 
        \[
        O\left(-1.5; -0.5\right)
        \]

    \end{enumerate}

    \item Démontrer qu'il s'agit d'un carré.

    Pour démontrer qu'il s'agit d'un carré, il faut vérifier que les côtés sont de même longueur et que les diagonales sont de même longueur et perpendiculaires.

    Calculons les longueurs des côtés et des diagonales :
    \[
    AB = \sqrt{(1 - (-2))^2 + (0 - 2)^2} = \sqrt{9 + 4} = \sqrt{13}
    \]
    \[
    BC = \sqrt{(-1 - 1)^2 + (-3 - 0)^2} = \sqrt{4 + 9} = \sqrt{13}
    \]
    \[
    CD = \sqrt{(-4 - (-1))^2 + (-1 - (-3))^2} = \sqrt{9 + 4} = \sqrt{13}
    \]
    \[
    DA = \sqrt{(-2 - (-4))^2 + (2 - (-1))^2} = \sqrt{4 + 9} = \sqrt{13}
    \]

    Les côtés sont tous égaux à $\sqrt{13}$.

    On peut également montrer que deux côtés adjascents seulement sont égaux. En effet, un parallélogramme
    a déjà ses côtés opposés deux à deux de même longueur.

    Calculons maintenant les longueurs des diagonales :
    \[
    AC = \sqrt{(-1 - (-2))^2 + (-3 - 2)^2} = \sqrt{1 + 25} = \sqrt{26}
    \]
    \[
    BD = \sqrt{(-4 - 1)^2 + (-1 - 0)^2} = \sqrt{25 + 1} = \sqrt{26}
    \]

    Les diagonales sont égales à $\sqrt{26}$.

    Donc $ABCD$ est un carré.

\end{enumerate}

\vspace{20pt}


% ========================================================
%             exercice 2
% ========================================================
\textbf{Exercice n$^\circ$2}

\begin{enumerate}

    \item Que peut-on dire des droites $(BC)$ et $(AD)$ ?

    Calculons les vecteurs $\overrightarrow{BC}$ et $\overrightarrow{AD}$ :
    \[
    \overrightarrow{BC} \Coord{-5 - 4}{1 - (-2)} = \overrightarrow{BC} \Coord{-9}{3}
    \]
    \[
    \overrightarrow{AD} \Coord{-1-2}{6-5} = \overrightarrow{AD} \Coord{-3}{1}
    \]

    On observe que $\overrightarrow{BC} = 3\overrightarrow{AD}$. Les vecteurs sont donc colinéaires, 
    ce qui signifie que les droites $(BC)$ et $(AD)$ sont parallèles.

    \vspace{5pt}
    On peut aussi calculer le déterminant et trouver 0.

    \item Déterminer les coordonnées des points $I$ et $K$.

    \begin{itemize}
        \item \textbf{Point I} :
        On sait que $\overrightarrow{IB} = -\dfrac{1}{2}\overrightarrow{BC}$.
        
        $\overrightarrow{BC}\Coord{-9}{3}$, donc $-\dfrac{1}{2}\overrightarrow{BC}
        \Coord{\dfrac{9}{2}}{-\dfrac{3}{2}}$

        Et $\overrightarrow{IB}\Coord{4-x_D}{-2-y_D}$

        On a donc :
        \[
        -\dfrac{1}{2}\overrightarrow{BC}
        \Coord{\dfrac{9}{2}}{-\dfrac{3}{2}} = \overrightarrow{IB}\Coord{4-x_D}{-2-y_D}
        \]
        
        En résolvant les équations sur les coordonnées, on trouve : $I\Coord{\dfrac{17}{2}}{-\dfrac{7}{2}}$

        \item \textbf{Point K} :
        On sait que $2\overrightarrow{BK} = \overrightarrow{BA} + \dfrac{1}{2}\overrightarrow{BC}$.
        
        On commence par écrire les coordonnées de $\overrightarrow{BK}\Coord{x_K-4}{y_K+2}$ 
        puis $2 \overrightarrow{BK}\Coord{2x_K-8}{2y_K + 4}$.
        
        Par ailleurs, en calculant $\overrightarrow{BA}$ puis $\dfrac{1}{2}\overrightarrow{BC}$ on trouve que
        $\overrightarrow{BA} + \dfrac{1}{2}\overrightarrow{BC}\Coord{-\dfrac{13}{2}}{\dfrac{17}{2}}$

        On pose alors les équations : 

        $2x_K-8 = -\dfrac{13}{2}$

        $y_K+2 = \dfrac{17}{2}$.

        Après avoir résolu on trouve : 

        $K\Coord{\dfrac{3}{4}}{\dfrac{9}{4}}$.
    \end{itemize}

    \item Démontrer que les points $I$, $K$ et $A$ sont alignés.

    Pour démontrer que les points $I$, $K$ et $A$ sont alignés, on vérifie que les vecteurs $\overrightarrow{IK}$ et $\overrightarrow{IA}$ sont colinéaires.
    
    Pour cela on utilise le calcul du déterminant : $det(\overrightarrow{IK};\overrightarrow{IA})$.

    \textbf{On peut également prendre d'autres vecteurs} comme $\overrightarrow{KA}$ et $\overrightarrow{AI}$ par exemple.

    Après résolution on trouve bien $0$.

    Donc les vecteurs sont bien colinéaires.

    Donc les points sont bien alignés.

    

\end{enumerate}

\vspace{20pt}


% ========================================================
%             Exercice 3
% ========================================================

\textbf{Exercice n$^\circ$3}

\vspace{5pt}
Cet exercice étant long, je ne mets ici que la méthode. À vous de faire les calculs associés.

\begin{enumerate}
    \item Déterminer les coordonnées : 
    \begin{enumerate}
        \item de S milieu de $[MP]$ : on utilise la formule du milieu 
        \item \textbf{faire un schéma pour s'aider!} R étant le symétrique, on a une égalité de vecteur qui peut s'écrire.
        On pose des inconnues et on résout.
        \item Ce genre de question nous laisse présumer qu'il s'agit sûrement d'un parallélogramme. 
        On fait un petit schéma et on écrit la propriété sur les égalités vectorielle.
        Si l'égalité est vérifiée c'est bien un parallélogramme.
    \end{enumerate}
    \item 
    \begin{enumerate}
        \item On utilise la colinéarité. Déterminant !
        \item On sait que $(MT)\ //\ (PR)$ car $PTMR$ est un parallélogramme. Donc si $M$, $T$ et $U$ sont 
        alignés, alors forcément $(UT)\ //\ (PR)$.
        \item On utilise la colinéarité.
        \item On utilise la formule sur les longueurs.
        \item On trouve que $TP = UR$. Or on sait aussi que $(UT)\ //\ (PR)$ et $(TR)\ //\ (UP)$. Donc le 
        quadrilatère est un parallélogramme qui a ses diagonales de même longueur : c'est un rectangle.
    \end{enumerate}
    \item 
    \begin{enumerate}
        \item On pose des inconnues, on écrit l'égalité vectorielle et on résout le système d'équation.
        \item Un point est une intersection à deux droites s'il est aligné avec les points appartenants à ces 
        deux droites. On montre donc que N, U et P sont alignés et que N, M et R sont alignés (avec la 
        colinéarité).
        \item On calcul les longueurs des côtés du triangle et on conclut (isocèle ou équilatéral en fonction 
        du résultat).
    \end{enumerate}
\end{enumerate}

\end{document}
